% from ega3-7-fr.tex
\oldpage[III]{189}
où $v'=\operatorname{Hom}({}^t d_p,1)$; lorsqu'on identifie canoniquement
$P_i\otimes M$ à $\operatorname{Hom}(\check P_i,M)$, $v'$ s'identifie donc à
$v=d_p\otimes1$, et l'on a par suite l'isomorphisme fonctoriel
\[
  T_p(M)\xrightarrow{\sim}\operatorname{Hom}_A(Q,M)
\]
cherché. En outre, $Q$, étant un quotient de $\check P_p$, est de type fini.
\end{proof}

\begin{proposition}[7.4.7]
Supposons vérifiées les conditions générales de (7.4.4). Alors, pour tout
$A$-module $M$ de type fini :
\begin{enumerate}
  \item[(i)] Les $T_i(M)$ sont des $A$-modules de type fini.
  \item[(ii)] Pour tout idéal $\mathfrak m$ de $A$, l'homomorphisme canonique
    \[
      \widehat{T_i(M)}
      \longrightarrow
      \varprojlim_n T_i\bigl(M\otimes_A(A/\mathfrak m^{n+1})\bigr)
      \tag{7.4.7.1}\label{III.7.4.7.1.fr}
    \]
    (où le premier membre est le séparé complété de $T_i(M)$ pour la topologie
    $\mathfrak m$-préadique) est bijectif.
\end{enumerate}
\end{proposition}

\begin{proof}
Comme dans (7.4.4), on se ramène d'abord au cas où les $P_i$ sont de type fini;
$A$ étant noethérien, les sous-modules des $P_i\otimes_A M$ sont de type fini,
d'où trivialement l'assertion (i). Quant à l'assertion (ii), elle résulte plus
généralement du lemme suivant :
\end{proof}

\begin{lemma}[7.4.7.2]
Soient $A$ un anneau noethérien, $u:E\to F$ un homomorphisme de $A$-modules de
type fini. Pour tout $A$-module de type fini, posons
\[
  K(M)=\operatorname{Ker}(u\otimes1_M),
  \qquad
  C(M)=\operatorname{Coker}(u\otimes1_M);
\]
alors les homomorphismes canoniques
\[
  \widehat{K(M)}\longrightarrow\varprojlim_n K(M_n),
  \qquad
  \widehat{C(M)}\longrightarrow\varprojlim_n C(M_n)
  \tag{7.4.7.3}\label{III.7.4.7.3.fr}
\]
(où l'on a posé
$M_n=M\otimes_A(A/\mathfrak m^{n+1})=M/\mathfrak m^{n+1}M$) sont bijectifs
pour tout idéal $\mathfrak m$ de $A$.
\end{lemma}

\begin{proof}
Comme $E\otimes M$ et $F\otimes M$ sont de type fini, et que le foncteur
$M\mapsto\widehat M$ est exact dans la catégorie des $A$-modules de type fini
($0_{\mathrm I}$, 7.3.3), $\widehat{K(M)}$ et $\widehat{C(M)}$ sont
respectivement le noyau et le conoyau de
\[
  \widehat{(u\otimes1)}:
  \widehat{(E\otimes M)}\longrightarrow\widehat{(F\otimes M)}.
\]
L'exactitude à gauche du foncteur $\varprojlim$ montre donc que
$\widehat{K(M)}=\varprojlim_n K(M_n)$; d'autre part, l'exactitude à droite du
produit tensoriel prouve que
$C(M_n)=C(M)\otimes_A(A/\mathfrak m^{n+1})$, donc
$\widehat{C(M)}=\varprojlim_n C(M_n)$ par définition.
\end{proof}

\begin{remark}[7.4.8]
Compte tenu de (6.10.5) et (6.10.6), on voit que, moyennant une hypothèse de
platitude supplémentaire, (7.4.7) redonne le fait que (4.1.7.1) est un
isomorphisme, c'est-à-dire l'essentiel du « premier théorème de comparaison »
pour les morphismes propres; en outre, l'énoncé s'applique non plus seulement
à un $\mathscr O_X$-Module cohérent, mais à un complexe de tels Modules. Il
serait intéressant d'obtenir un énoncé comprenant à la fois (7.4.7) et
(4.1.7.1) comme cas particuliers. On notera que lorsque les $P_i$ sont de type
fini, la démonstration de (7.4.7) n'utilise pas le fait que ce sont des modules
plats; il y aurait lieu d'examiner si la conclusion de (7.4.7) est encore
valable lorsque les $P_i$ ne sont pas supposés plats ni de type fini, mais que
les $H_i(P_\bullet)$ sont supposés de type fini pour tout $i$ et nuls pour
$i<i_0$. Est-il alors possible de remplacer $P_\bullet$ par un complexe
$P'_\bullet$ de $A$-modules de type fini tel que les foncteurs
$H_\bullet(P_\bullet\otimes M)$ et $H_\bullet(P'_\bullet\otimes M)$ (qui ne
sont plus homologiques) soient encore isomorphes ?
\end{remark}

